\begin{pastexam}[label={px:ccu-113-csb-05}]{中正資工所乙組 113 學年度 第 9 題}
（單選，並說明理由）

(5\%) What will we get after heapifying the integer array \{30, 10, 45, 25, 60, 15, 35, 5\} to a Max
Heap? Note that we use the bottom-up method that runs in $O(n)$ time.

(a)~\{5, 10, 15, 25, 30, 35, 45, 60\} \qquad (b)~\{60, 45, 35, 30, 25, 15, 10, 5\}\\
(c)~\{60, 45, 35, 10, 25, 15, 30, 5\} \qquad (d)~\{60, 30, 45, 25, 10, 15, 35, 5\}
\end{pastexam}

\begin{pastsoln}
答案是 (d)。

由下而上建堆積（資結書虛擬碼~8.3）：陣列 $A[1..8]$ 看成完整二元樹，$A[i]$ 的子節點是 $A[2i]$、
$A[2i+1]$。 從最後一個有子節點的位置 $\lfloor 8/2 \rfloor = 4$ 往回到 $1$，每個位置做一次下沉
（和比較大的子節點比，比它小就交換，一直往下）。
\begin{center}
\small
\begin{tabular}{clc}
\toprule
$i$ & 動作 & 之後的陣列 \\
\midrule
（開始） & & $30, 10, 45, 25, 60, 15, 35, 5$ \\
4 & $A[4] = 25$，子節點 $A[8] = 5$，不動 & $30, 10, 45, 25, 60, 15, 35, 5$ \\
3 & $A[3] = 45$，子節點 $15, 35$，不動 & $30, 10, 45, 25, 60, 15, 35, 5$ \\
2 & $A[2] = 10$，子節點 $25, 60$，和 $60$ 換；$A[5]$ 沒有子節點 & $30, 60, 45, 25, 10, 15, 35, 5$ \\
1 & $A[1] = 30$，子節點 $60, 45$，和 $60$ 換； & \\
  & 到 $A[2] = 30$，子節點 $25, 10$，停 & $60, 30, 45, 25, 10, 15, 35, 5$ \\
\bottomrule
\end{tabular}
\end{center}
結果 $\{60, 30, 45, 25, 10, 15, 35, 5\}$，就是 (d)。 檢查：$60 \ge 30, 45$；$30 \ge 25, 10$；
$45 \ge 15, 35$；$25 \ge 5$，每個父節點都不小於子節點。

其他選項：
\begin{itemize}
\item (a) 是由小到大排好的陣列，根是最小值，不是最大堆積。
\item (b) 由大到小排好的陣列確實是一個最大堆積，但那是排序的結果；建堆積只做局部的下沉，
不會把整個陣列排好（建堆積只要 $O(n)$，排序需要 $\Omega(n \lg n)$ 次比較）。
\item (c) 也是一個合法的最大堆積，但它是\textbf{一個一個插入}（每次上浮）得到的：依序插入
$30, 10, 45, 25, 60, 15, 35, 5$，結果正是 $\{60, 45, 35, 10, 25, 15, 30, 5\}$。 題目指定的是
$O(n)$ 的由下而上建法，兩種建法的結果不一定相同。
\end{itemize}
\end{pastsoln}
